\chapter{Deployment}
\label{ch:deployment}
\section{Hosting}
The system is deployed as a Docker container on Render's free tier (750 hours/month, 512MB RAM, no credit card required, roughly 30–60 second cold-start after 15 minutes of inactivity). This platform was chosen only after a genuinely thorough comparison, redone twice as circumstances changed: Railway, Fly.io, Hugging Face Spaces, Back4App, and Google Cloud Run were each evaluated and ruled out for concrete reasons (a gutted free-credit model, a mandatory credit card, Docker support behind a paywall, insufficient free RAM, and mandatory account/billing setup, respectively); Vercel, Netlify, GitHub Pages, and Cloudflare were separately ruled out because they are serverless/static platforms that cannot run this project's Dockerized, stateful backend at all; and Koyeb, initially a strong secondary candidate, was dropped after its mid-acquisition by Mistral AI was confirmed via its own public announcement — judged too great a stability risk to build on immediately before a graded deadline.

The first real deployment attempt failed with an out-of-memory kill on Render's 512MB limit. The root cause was that the backend re-parsed the entire 79-page PDF from scratch (a roughly 12.7-second, memory-heavy operation) on every process startup. The fix moves this work out of the request path entirely: a pre-built, validated node cache (\texttt{data/\allowbreak{}processed/\allowbreak{}nodes.json}) is baked into the Docker image at build time, and the backend loads that cache at startup, falling back to a live parse only if no cache is present — verified to be perfectly lossless (zero mismatches across all 327 nodes) before being trusted.

Because the local Ollama option requires a locally-running model with real RAM that no free hosting tier provides, the hosted deployment runs Groq-only; the local option remains fully available and tested for anyone running the project on their own machine via Docker Compose.

\section{Containerization}
The project ships as a two-stage Docker build: a \texttt{node:20-slim} stage builds the React frontend, and the final \texttt{python:3.11-slim} stage contains only the built static assets and the Python backend — Node itself never ships in the final image. \texttt{docker-compose.yml} defines an \texttt{app} service and, for local development, an \texttt{ollama} service on a shared network; the Groq option needs no local container at all. Configuration is entirely environment-variable driven (\texttt{.env}, excluded from both git and the Docker build context) with a documented \texttt{.env.example} template, so no API key is ever at risk of being committed or baked into an image layer.

\section{Web application}
A thin FastAPI backend (\texttt{webapp/\allowbreak{}backend.py}) wraps the agent's \texttt{answer\_\allowbreak{}question()} function with no independent business logic of its own — every real decision already lives in the modeling and agent layers described in Section 8. The frontend was rewritten partway through the project from static HTML/CSS/JS to React (via Vite) plus GSAP animation, an explicit trade-off the intern chose after being shown both options, accepting a new build-toolchain dependency in exchange for a more maintainable, componentized codebase. It ships as three independent single-page applications (landing, chat, dashboard) rather than one combined app, matching the backend's existing separate routes for each page.

The chat interface supports full multi-conversation history (a sidebar of past conversations persisted in the browser, in the style of a typical AI chat product), a live model-status indicator, a picker for which LLM mode to use, and a toggle to compare the LLM-phrased answer against the original deterministic template answer side by side — a deliberate transparency feature, not a debugging leftover, letting a user see exactly what changed between the verified facts and their conversational rephrasing. A companion analytics dashboard visualizes the DAM's own structure (node-type breakdown, authority-code distribution, most-referenced roles, unresolved-role rate) with a chapter filter and linked charts, built specifically to showcase the underlying data model rather than agent usage statistics.

\section{Multi-language support}
Following direct feedback that non-English speakers could not use the system, the agent now supports French, Spanish, Portuguese, and Arabic alongside English, via an explicit language picker in the chat interface (in addition to automatic detection of the question's own language). The architecture deliberately does not attempt to rebuild the underlying retrieval pipeline per language — doing so would mean re-validating years of accumulated, hard-won English-language parsing logic in four more languages with no additional labeled data. Instead, a non-English question is translated to English before it reaches retrieval (which remains completely unchanged and unaffected), and the final answer is translated into the requested language afterward. The grounding check described in Section 8.4 is extended, not weakened, for this case: it still requires the literal English role names to appear verbatim in the model's output regardless of what language the surrounding sentence is written in, and any answer that fails this check falls back to the safe English template rather than showing an unverified translation. Arabic additionally triggers right-to-left layout in the chat interface. A small, fixed set of interface labels (roughly 15 strings) is hand-translated once into each language rather than translated live by the LLM on every use — deliberately, for speed, cost, and consistency, following the same reasoning applied elsewhere in the project to small, fixed vocabularies (Section 8.6's tone-response prefixes use the identical pattern).

\section{Operational incident: a live production bug and its resolution}
In September 2026 the deployed instance began returning HTTP 404 errors from the Groq API on every LLM-mode request. Root-causing this against Groq's own current documentation (rather than assumed prior knowledge, since API model availability changes over time) confirmed that the model in use, \texttt{llama-3.3-70b-versatile}, had been moved to enterprise-only pricing on 16 August 2026 and was no longer available on the free tier being used. The fix — migrating the default model to \texttt{openai/\allowbreak{}gpt-oss-120b}, Groq's own documented free-tier replacement — was applied across the code, environment configuration, documentation, and tests. This incident is a useful, honest illustration of a real operational risk in any system that depends on a third-party model API: model availability and pricing tiers are not static, and a system built against one can silently break with no code change of its own.

\section{Chatbot feature iteration following professor and user feedback}
Following review by the internship supervisor, four further improvements were requested and delivered, each following the same discipline of confirming a real gap before building the fix:

\begin{itemize}
\item \textbf{Vague-question and out-of-scope detection.} A question too short or generic to confidently resolve (measured, for example, by content-word count after removing stopwords and by the score gap between the top match and its nearest competitor) now returns a clarification request listing real candidate activities rather than a low-confidence guess; a question genuinely unrelated to the DAM's subject matter is reported as out of scope rather than matched to an unrelated activity by incidental word overlap. Both are explicitly disclosed as best-effort heuristics, not true topic understanding — a small number of edge cases (e.g. "tell me a joke") are known to be misclassified, but the system's core honesty guarantee (never fabricating a false DAM answer) holds regardless.
\item \textbf{Mandatory Check/Verify and informed-party notes.} Following an explicit domain-knowledge correction from the intern (drawn from direct internship experience, not derivable from the PDF's formatting alone) — that a Check/Verify entity on an activity's row is mandatory information, not optional supplementary detail — every intent-specific answer now automatically appends a note naming the mandatory checking/verifying role and, when present and relevant, the informed parties, regardless of which specific authority action was originally asked about. This was scoped, after discussion, to apply to every intent uniformly (a property of the activity itself, not of the question asked) and to be applied deterministically and consistently rather than only "when it seems useful" — chosen specifically because a consistent, testable rule is more trustworthy than a heuristic that might silently omit the note on an unusual phrasing.
\item \textbf{Action-code legend explanation.} A question asking what the authority-code letters themselves mean (e.g. "what's I, A and (i)?") was, for a period, incorrectly swallowed by the conversational-context-carryover logic (a short question containing only single letters looked, by an internal word-count heuristic, indistinguishable from a genuinely vague follow-up). The fix adds a dedicated detection path that runs before context-carryover is ever considered, reusing the DAM's own already-extracted authority-code reference data rather than a separately hand-maintained explanation.
\item \textbf{Emotional tone detection.} Every response, regardless of whether it is a factual DAM lookup, a glossary answer, or an honest refusal, is now checked against a cheap, deterministic pre-filter for signs of user frustration or confusion (repeated punctuation, all-caps words, characteristic phrases); only messages that pass this filter trigger a single, cheap LLM classification call, and only that call's result (never free-form LLM text) determines whether a short, deterministic, language-matched empathetic opener is prepended to the answer. This keeps the added warmth completely outside the factual grounding-check boundary — it can never itself introduce an ungrounded claim, by construction.
\end{itemize}
\section{Live status and reliability signaling}
The web interface polls the backend's health endpoint continuously and shows a real-time status indicator (rather than assuming the backend is always reachable), and every answer's metadata (which LLM mode actually produced it, whether translation succeeded, whether the answer fell back to the deterministic template) is surfaced directly in the UI rather than hidden — a deliberate transparency choice consistent with the project's overall stance that an honest "this didn't fully work" is always preferable to a confident but unverifiable answer.

\section{The hosting-platform search, narrated in full}
Section 10.1 summarizes the final hosting decision; the process of reaching it is worth narrating because it was genuinely re-done twice as circumstances changed, not settled once and left alone. The first pass ruled out Railway (a free tier gutted to roughly one dollar of monthly credit, no longer practically free), Fly.io (requires a credit card even for its smallest tier), Hugging Face Spaces (Docker support is a paid feature), Google Cloud Run (requires a card and more account setup than the project's scope warranted), and Back4App (256MB of RAM, judged too tight for this application's needs). Render was selected from what remained.

A second pass was triggered later by further research into alternatives before the first Render deployment attempt had even been retried after its initial out-of-memory failure (Section 10.1). This pass considered Vercel (a serverless-functions platform whose "FastAPI preset" ignores a repository's own Dockerfile entirely, meaning neither the frontend build step nor the node-cache build step this project depends on would ever run, and whose cold-start-per-invocation model has no equivalent for this application's single in-memory server state), GitHub Pages (static files only, by design incapable of running any backend code at all), Netlify (the same serverless/Dockerfile mismatch as Vercel), and Cloudflare (whose Containers product would actually fit, but sits behind a five-dollar-a-month paid plan with no free equivalent). Koyeb was, for a period, considered the strongest secondary candidate — architecturally the same Docker-native model as Render — until its dashboard was found to display a banner announcing its acquisition by Mistral AI and a mid-transition "stay tuned for a revamped experience" state instead of a normal service-creation flow; judged too great a stability risk to build a graded deliverable on top of during an active platform transition, this was set aside rather than pursued further. The search concluded, a second time, back at Render — the platform had been correct on the first pass, and the out-of-memory failure that had briefly cast doubt on it was a fixable bug in this project's own startup behavior (Section 10.1), not a genuine limitation of the platform itself.

\section{Glossary and abbreviation lookup}
Following direct testing feedback that questions like "what does DDG mean" returned a generic failure, the project extracted the DAM's own Abbreviations list (document pages 2–7, deliberately scoped separately from the longer, line-wrapping Glossary section on pages 8–11, which was set aside as a harder extraction problem not worth the time budget at that stage) into a structured reference file of 170 entries, and built a dedicated glossary-lookup path (\texttt{agent/\allowbreak{}glossary.py}) alongside inline acronym expansion within role names elsewhere in the system's answers — 19 of the DAM's 131 distinct canonical role names are themselves entries in this same abbreviations list. Two real parsing bugs were found and fixed while building the extractor: rows containing only a definition continuation (no new acronym on that line) were incorrectly being treated as starting a new entry, and page-footer roman numerals were leaking into the tail of a definition's text. A same-day follow-up round of direct testing found the glossary trigger phrase "what does X stand for" did not recognize the common variant "what does X stands for" (a minor grammatical slip a real user actually typed), and separately that "RDNG" — an acronym that appears throughout the DAM's own role names but, unlike "RDG", is never separately defined in the DAM's own printed Abbreviations pages — could not be resolved at all. Both were fixed: the trigger regex was widened, and a small, explicitly separate \texttt{MANUAL\_\allowbreak{}ALIASES} dictionary was introduced for the handful of acronyms like this one, kept deliberately apart from the file that stores only what the DAM document itself actually defines, so the system's own added knowledge is never silently mixed with, or mistaken for, verified source-document content.

\section{The React frontend rewrite, and a real sandbox environment limitation}
The original frontend was static, build-free HTML/CSS/JavaScript. After showing both options and their trade-offs, the project was rewritten to React (via Vite) with GSAP animation — accepting a new Node-based build toolchain in exchange for a more maintainable, componentized codebase, an explicit trade-off decision rather than a default choice. The rewrite ships as three independent single-page applications (landing, chat, dashboard), matching the backend's existing separate routes for each page rather than consolidating into one combined app.

Building this inside the development sandbox surfaced a genuine, non-obvious environment problem: running \texttt{npm install} inside a folder synchronized by OneDrive was unworkably slow, because the cloud-sync layer intercepts and processes every one of the thousands of small files \texttt{node\_\allowbreak{}modules} produces. The practical workaround, used for every subsequent frontend build in this project, was to run the install and build steps in an unsynchronized temporary directory and copy back only the finished, built output — never the intermediate \texttt{node\_\allowbreak{}modules} tree itself. A second, related limitation was that browser-based UI testing (via Playwright) was attempted but blocked by the sandbox's outbound network allowlist, which prevented downloading a Chromium binary; the project's automated test coverage for the frontend was consequently redesigned around source-level structural checks (\texttt{tests/\allowbreak{}test\_\allowbreak{}frontend\_\allowbreak{}source.py}) rather than rendered-browser assertions, an honest adaptation to a real tooling constraint rather than an attempt to work around it invisibly.

Direct testing of the freshly rebuilt frontend surfaced five further real, user-visible gaps in a single round, each confirmed against a real screenshot before being fixed: a widened glossary-trigger regex was still missing the specific phrasing "what does RDG stands for" on the new build; "RDNG" needed the same manual-alias treatment as above; the chat page rendered its message column at the full page width with an ungainly gray gap on wide screens, fixed with an explicit centered, max-width inner container; there was no way to see or return to a previous conversation, addressed by building a full multi-conversation sidebar backed by browser local storage; and the agent's replies, while factually sound, read as flat and mechanical, addressed with varied smalltalk phrasing and a warmer (but still fact-grounded) system prompt for LLM-phrased answers.

\section{Deterministic typo tolerance}
Real testing surfaced a cluster of ordinary human typing mistakes the system could not handle: a misspelled "aproves", a misspelled "qaurterly mision", a truncated "helo", and a garbled "waht does DDG men" — each causing an otherwise-correct question to fail to resolve. Rather than route typo correction through the language model (which would make even this basic layer of the system dependent on network access and cost), a deterministic correction utility (\texttt{knowledge/\allowbreak{}typo\_\allowbreak{}correct.py}) was built on Python's standard-library \texttt{difflib} similarity scoring, with two different similarity thresholds calibrated for two different situations: a looser threshold for small, curated vocabularies (smalltalk trigger words, glossary trigger words, intent-detection keywords), and a stricter threshold for the large, organic vocabulary of real DAM title words used in text search, where a looser threshold risks false corrections. That false-correction risk was not hypothetical: an early threshold setting was found to incorrectly "correct" the genuinely different word "unrelated" into "related" due to their high surface similarity, caught and fixed by raising the large-vocabulary threshold until real typo cases remained correctly matched while this specific false positive no longer was. Words shorter than four letters and all-capitalized words (likely to be genuine acronyms rather than misspellings) are deliberately excluded from correction, a conscious, documented trade-off rather than an oversight.

\section{Conversational context carryover: a two-round fix}
Section 8.6 references this fix; the two rounds it took to get right are worth narrating because the first, reasonable-seeming fix was defeated by a real follow-up message within the same testing session, which is itself a useful illustration of why every fix in this project was tested against real, adversarial-feeling input rather than a single happy-path example. Direct testing found that after resolving one activity and asking a follow-up question referencing "that activity" without repeating its name, the system resolved the follow-up to a completely different, unrelated activity that merely happened to share enough surface vocabulary to score above zero in text search. The first fix recognized a small set of explicit anaphoric trigger phrases ("that activity", "it", "this") and carried the previous activity's identity forward whenever one was detected. This fix was defeated on the very next real message: a follow-up phrased with a typo and no recognized trigger phrase ("and who are the informed partie?") landed on the exact same wrong, unrelated activity at an almost identical confidence score to the original failure — proof that the confidence score itself, not the specific wording of the follow-up, was the real signal that should have been used from the start. The second, final fix discarded the phrase-list approach entirely and instead always resolves the fresh question first, only overriding with the previous turn's activity when the fresh resolution's own confidence score falls below a threshold calibrated against real measured examples of genuine matches versus real measured examples of coincidental low-confidence mismatches. This is explicitly documented as a sharper but still imperfect trade-off: a genuinely new question that happens to be weakly worded can still be swept into carryover, accepted because such a question would have been a low-confidence, borderline answer on its own regardless.

\section{Mobile responsiveness and the dashboard rework}
Following the React rewrite, the interface was adapted for small mobile viewports: a hamburger-triggered off-canvas conversation sidebar below a defined breakpoint, a responsive header that wraps its controls rather than overflowing, and message bubbles that use a percentage-based width on mobile instead of the fixed desktop width. Separately, following direct feedback from the internship supervisor requesting a more dynamic, interactive, and clearly relevant dashboard, the analytics page was substantially reworked: a chapter filter re-scopes every card and chart to a single chapter of the DAM (graph-level totals remain whole-document, since cross-chapter reference edges cannot be meaningfully attributed to one chapter alone), two new coverage metrics were added (average responsibilities per node, and the percentage of nodes with no directly recorded responsibilities — a genuine, previously invisible data-quality signal about the extraction itself, not a cosmetic addition), and clicking a bar on the authority-code distribution chart now filters the roles chart above it to show who is most active under that specific code. Building this rework caught one further real bug: a chart-rendering component had an empty effect-dependency array, meaning its charts silently never updated after their very first render regardless of new data — fixed by correcting the dependency list to include the actual data and options the chart depends on.

